\documentclass[10pt]{article}
\usepackage[T1]{fontenc}
\usepackage[utf8]{inputenc}
\usepackage[french]{babel}
\usepackage{amsmath,amssymb,mathrsfs,mathpazo}
\usepackage{geometry}
\geometry{a4paper,margin=2cm}
\usepackage{graphicx,array,tabularx,tabularray,colortbl,multirow,multicol,booktabs,enumitem,pifont,ulem,fancybox,diagbox,cancel}
\usepackage{pstricks,pstricks-add,pst-plot,pst-tree,pst-node,pst-text,pst-func,pst-math,pst-bspline,pst-3dplot,pst-eucl,pst-solides3d}
\let\pscalculate\undefined
\usepackage{pst-bezier}
\usepackage{tikz,pgfplots,tkz-tab}
\usetikzlibrary{arrows,arrows.meta,patterns,calc,positioning,trees,decorations.markings,shapes.geometric,shapes.misc,shapes.arrows,intersections,decorations.pathmorphing,decorations.pathreplacing,backgrounds,shadows,babel}
\pgfplotsset{compat=1.18}
\usepgfplotslibrary{fillbetween}
\usepackage[np]{numprint}
\usepackage{listings,makecell,pas-tableur,twoopt}
\let\barre\undefined
\usepackage[pstricks]{tablvar}
\definecolor{darkgreen}{RGB}{0,100,0}
\providecommand{\red}{\color{red}}
\providecommand{\blue}{\color{blue}}
\providecommand{\green}{\color{green}}
\providecommand{\black}{\color{black}}
\providecommand{\white}{\color{white}}
\providecommand{\gray}{\color{gray}}
\providecommand{\cyan}{\color{cyan}}
\providecommand{\magenta}{\color{magenta}}
\providecommand{\yellow}{\color{yellow}}
\providecommand{\violet}{\color{violet}}
\providecommand{\orange}{\color{orange}}
\providecommand{\euro}{\texteuro}
\providecommand{\EUR}{\texteuro}
\providecommand{\ieme}{\textsuperscript{e}}
\providecommand{\degres}{\ensuremath{^\circ}}
\providecommand{\degree}{\ensuremath{^\circ}}
\providecommand{\vv}[1]{\overrightarrow{#1}}
\providecommand{\Coord}[1]{#1}
\providecommand{\decosix}{}
\providecommand{\starredbullet}{\ensuremath{\bullet}}
\pagestyle{empty}
\setlength{\parindent}{0pt}
\usepackage{adjustbox}\DeclareUnicodeCharacter{2236}{:}\DeclareUnicodeCharacter{2190}{\ensuremath{\leftarrow}}
\providecommand{\R}{\mathbb{R}}
\providecommand{\N}{\mathbb{N}}
\providecommand{\Z}{\mathbb{Z}}
\providecommand{\Q}{\mathbb{Q}}
\providecommand{\D}{\mathbb{D}}
\providecommand{\C}{\mathbb{C}}
\providecommand{\e}{{\mathrm{e}}}
\providecommand{\vect}[1]{\overrightarrow{#1}}
\providecommand{\vectt}[1]{\overrightarrow{\mathrm{#1}}}
\providecommand{\vvt}[1]{\overrightarrow{\mathrm{#1}}}
\providecommand{\np}[2][]{{#2~#1}}
\providecommand{\vv}[1]{{\overrightarrow{#1}}}
\renewcommand{\d}{{\mathrm{d}}}
\providecommand{\degres}{{^\circ}}
\providecommand{\degre}{{^\circ}}
\providecommand{\blue}{\color{blue}}
\providecommand{\red}{\color{red}}
\providecommand{\enskip}{\quad}
\providecommand{\llbracket}{\lbrack\!\lbrack}
\providecommand{\rrbracket}{\rbrack\!\rbrack}
\providecommand{\euro}{{\text{€}}}
\providecommand{\no}{{\text{n°}}}
\providecommand{\interval}[2]{{[#1;#2]}}
\providecommand{\intervFO}[2]{{[#1;#2[}}
\providecommand{\dotfill}{\ldots}

\providecommand{\ds}{\displaystyle}
\providecommand{\textding}[1]{#1}
\providecommand{\Oij}{\(\mathrm{O};\vec{\imath},\vec{\jmath}\)}
\providecommand{\Oijk}{\(\mathrm{O};\vec{\imath},\vec{\jmath},\vec{k}\)}
\providecommand{\Ouv}{\(\mathrm{O};\vec{u},\vec{v}\)}
\providecommand{\intervFF}[2]{\ensuremath{[#1;#2]}}
\renewcommand{\euro}{\ifmmode\text{\texteuro}\else\texteuro\fi}
\DeclareUnicodeCharacter{00B0}{\ifmmode^{\circ}\else\textdegree\fi}
\DeclareUnicodeCharacter{20AC}{\ifmmode\text{\texteuro}\else\texteuro\fi}
% Énoncés : Éducation nationale. Transcription et corrigés : APMEP, Denis Vergès et contributeurs. Découpage et classement : CoopMaths, Sébastien Lozano, Éric Elter et contributeurs. Adaptation de la présentation : Maths Explorer.
% CC BY-SA : https://coopmaths.fr/www/about/
% https://coopmaths.fr/alea/static/e3c/2020/tex/e3c_2020_00_sujet64_3_cor.tex
\begin{document}
\noindent\textit{APMEP / CoopMaths. CC BY-SA. Présentation adaptée par Maths Explorer.}\par\medskip
\subsection*{1.}

On a :
\begin{align*}
&M(x \,;\, y) \in \mathcal{C} \\
\iff &AM^2 = 5^2 \\
\iff &(x - 2)^2 + (y - 5)^2 = 25 \\
\iff &x^2 - 4x + 4 + y^2 - 10y + 25 = 25 \\
\iff &x^2 + y^2 - 4x - 10y = -4.
\end{align*}

\subsection*{2.}

\begin{align*}
&B(5 \,;\, 9) \in \mathcal{C} \\
\iff &5^2 + 9^2 - 4 \times 5 - 10 \times 9 = -4 \\
\iff &25 + 81 - 20 - 90 = -4 \\
\iff &106 - 110 = -4,
\end{align*}
qui est vrai.

\subsection*{3.}

A est le centre du cercle et B est un point de ce cercle. On sait que la tangente en un point du cercle est perpendiculaire au rayon contenant ce point \(B\).

\subsection*{4.}

Si \(\mathcal{T}_B\) est cette tangente, on a :
\[
M(x \,;\, y) \in \mathcal{T}_B \iff (\overrightarrow{BM}) \perp (\overrightarrow{AB}) \iff \overrightarrow{BM} \cdot \overrightarrow{AB}.
\]
Avec \(\overrightarrow{BM} \begin{pmatrix} x - 5 \\ y - 9 \end{pmatrix}\) et \(\overrightarrow{AB} \begin{pmatrix} 3 \\ 4 \end{pmatrix}\), on a donc :
\begin{align*}
&\overrightarrow{BM} \cdot \overrightarrow{AB} = 0 \\
\iff &3(x - 5) + 4(y - 9) = 0 \\
\iff &3x + 4y - 15 - 36 = 0 \\
\iff &3x + 4y - 51 = 0.
\end{align*}

\subsection*{5.}

Un point de l'axe des ordonnées est caractérisé par son abscisse nulle (\(x = 0\)), donc :
\[
\begin{cases}
x^2 + y^2 - 4x - 10y = -4 \\
x = 0
\end{cases}
\Rightarrow y^2 - 10y = -4
\iff y^2 - 10y + 4 = 0.
\]
Pour cette équation :
\[
\Delta = 100 - 4 \times 4 = 84 = 4 \times 21 = (2\sqrt{21})^2 > 0,
\]
il y a donc deux racines :
\[
y_1 = \frac{10 + 2\sqrt{21}}{2} = 5 + \sqrt{21} \quad \text{et} \quad y_2 = \frac{10 - 2\sqrt{21}}{2} = 5 - \sqrt{21}.
\]
Le cercle \(\mathcal{C}\) a deux points communs avec l'axe des ordonnées de coordonnées :
\[
(5 + \sqrt{21} \,;\, 0) \quad \text{et} \quad (5 - \sqrt{21} \,;\, 0).
\]

\end{document}
